% ===== 原文 pp.21–22（PDF pp.22–23）=====
% tr-02 片段 c：(4.16)–(4.20)、4.3 相（Calabi–Yau情形）
% [原文 p.21]
其中$W_\zeta$是$\mu^{-1}(\zeta)/G$上由$W$诱导的函数。此外，若$W_\zeta$是Bott–Morse的，则理论进一步约化为$\Crit(W_\zeta)$上的非线性sigma模型，而$\Crit(W_\zeta)$恰好就是Higgs分支$X_\zeta$。这样的相称为几何相。在腔室之间的壁（相边界）上，$\Crit(W)\cap\mu^{-1}(\zeta)$的某些位置处存在连续的稳定子群。于是经典真空空间中会发展出一个分量，称为Coulomb分支，在其中$\sigma$可以取稳定子群的嘉当子代数中的任意值。这个非紧空间的出现可以视为一种奇异性。

量子效应会对此图像作出显著修正。特别地，经典Coulomb分支可能被抬升；或者即使没有经典Coulomb分支，也可能出现量子Coulomb真空。为看清这一点，我们考察场空间中$\sigma$在$\fg_\C$的嘉当子代数$\ft_\C$中取大的、一般的值的区域。此时 (4.14) 的第二、三项使$\phi$的许多分量（通常是全部分量，我们暂作此假设）获得质量。积分掉有质量模式后，我们得到仅由极大环面$T$的规范多重态构成的有效理论，其有效扭超势为
\begin{equation}
\widetilde{W}_{\text{eff}}(\sigma)=-t(\sigma)+2\pi i\rho(\sigma)-\sum_iQ_i(\sigma)\left(\log(Q_i(\sigma)/\mu)-1\right).
\tag{4.16}
\end{equation}
这里$\rho=\frac{1}{2}\sum_{\alpha>0}\alpha$是$\fg$的正根之和的一半，$Q_i$是$V$的权，而$t$是标度$\mu$处的FI-theta参数。有效理论的规范耦合常数是$\sigma$的一个复杂函数$e_{\text{eff}}(\sigma)$，但当$|\sigma|$很大时它趋于给定值$e$。有效势为
\begin{equation}
U_{\text{eff}}(\sigma)=\min_{n\in\mathrm{P}}\frac{e^2_{\text{eff}}(\sigma)}{2}\left|\dd\widetilde{W}_{\text{eff}}(\sigma)+2\pi in\right|^2
\tag{4.17}
\end{equation}
注意，(4.16) 中对数分支的选取没有物理效应——不同的选取会把$\widetilde{W}_{\text{eff}}(\sigma)$移动$2\pi i\mathrm{P}(\sigma)$中的一个元素，但这不影响 (4.17)。当$U_{\text{eff}}(\sigma)=0$（即$\dd\widetilde{W}_{\text{eff}}(\sigma)\in2\pi i\mathrm{P}$）时，点$\sigma$才是量子理论中真正的Coulomb真空。特别地，经典Coulomb分支可能在量子理论中不复存在；或者即使没有经典Coulomb分支，也可能出现真正的Coulomb真空。为完整起见，还应当考察场空间中$\sigma$在某个子群$H\subset G$的嘉当子代数中取大的、一般的值、同时$\phi$的$H$-中性分量也取大值的区域，并检查是否存在真正的混合Coulomb-Higgs真空。

% [Tutor Note: p.21 式 (4.16)–(4.17) 中“对数分支无关”的判据：把$\log$的宗量改乘$\ee^{2\pi ik_i}$（$k_i\in\Z$），第$i$项移动$-2\pi ik_iQ_i(\sigma)$，这是$2\pi i\mathrm{P}(\sigma)$中的元素；而 (4.17) 里的$\min_{n\in\mathrm{P}}$恰好允许$\dd\widetilde{W}_{\text{eff}}$作同样的$2\pi i\mathrm{P}$平移，于是$U_{\text{eff}}$逐点不变。核对方法：把$n\to n-\sum_ik_iQ_i$代换后 (4.17) 的方括号不变。]

理论的性质极大地依赖于Calabi–Yau条件的无穷小版本$\fg\subset\fsl(V)$是否满足，也就是说，依赖于由
\begin{equation}
b_1(X):=\tr_V(X)\qquad X\in\fg
\tag{4.18}
\end{equation}
定义的$b_1$是否为零。

% [原文 p.22]
注意，$b_1$可视为$\mathrm{P}^W\cong\mathrm{P}_{\mathrm{Z}_G}$中的一个元素。作为$\mathrm{P}^W$的元素，它也可以写成$b_1=\sum_iQ_i$。

% [Tutor Note: p.22 的$b_1=\sum_iQ_i$可用来核对 (4.18)：$b_1(X)=\tr_V(X)$在$G$的共轭作用下不变，故落在Weyl不变子空间$\mathrm{P}^W$中；把$V$按权分解，$\tr_V(X)=\sum_iQ_i(X)$恰是权重之和，这正是“$\mathrm{P}^W\cong\mathrm{P}_{\mathrm{Z}_G}$中的一个元素”的显式形式。核对方法：取$G=U(1)$、$V=\C$权为$Q$，则$b_1=Q$；取$V=\C^N$标准表示则$b_1=0$，对应$SL(V)$条件。]

\para{Calabi–Yau 情形}

假定它满足，即$b_1=0$。此时FI参数在重整化群下不变，并且量子理论中存在轴$U(1)$ R-对称性。相应地，(4.16) 中的$\widetilde{W}_{\text{eff}}$与参数$\mu$无关。特别地，真空方程$\dd\widetilde{W}_{\text{eff}}(\sigma)\in2\pi i\mathrm{P}$在标度变换$\sigma\to\lambda\sigma$（$\lambda\in\C^\times$）下不变。这意味着，若$\sigma$是一个Coulomb真空，则它的任何伸缩也仍是Coulomb真空。特别地，这类真空构成的空间即Coulomb分支，必须是非紧的。此外，Coulomb分支的存在对FI-theta参数$t$施加一个非平凡的约束。事实上，方程$\dd\widetilde{W}_{\text{eff}}(\sigma)\in2\pi i\mathrm{P}$给出用$\sigma$的坐标之比对$t$的参数表示，从而在$t$的空间中定义一个复超曲面。设$\Delta\subset\fg^{*G}_\C/2\pi i\Lambda_G$为判别轨迹，在其上存在Coulomb分支和/或混合分支。它是若干超曲面的并。当投影到$\zeta$空间时，判别轨迹$\Delta$投影为一个变形虫（amoeba），它渐近地趋于相边界。因此，$t$的正则值空间为
\begin{equation}
\mathfrak{M}_t=\fg^{*G}_\C/2\pi i\Lambda_G-\Delta.
\tag{4.19}
\end{equation}

由于$\Delta\subset\fg^{*G}_\C/2\pi i\Lambda_G$具有复余维一，人们可以从一个相走到另一个相而不碰到它。特别地，不同的相之间没有尖锐的相变。

该理论在红外极限下流到一个SCFT，其$\hat c=\tr_V(1-R)-\dim G$，而FI-theta参数（相应地，$W$的参数）正是该SCFT的精确边缘A-项（相应地，B-项）形变参数。也就是说，它们参数化了SCFT模空间的子流形
\begin{equation}
\mathfrak{M}_t\subset\mathfrak{M}^0_A,\qquad\mathfrak{M}_W\subset\mathfrak{M}^0_B.
\tag{4.20}
\end{equation}
往往，这里的包含关系$\subset$就是等号$=$。

在一个几何相中，$\mathfrak{M}_t$与$\mathfrak{M}_W$分别是Higgs分支的复化Kähler类模空间与复结构模空间（的一部分）。

% [Tutor Note: 本片段范围内原文的拼写笔误（正文照录，不代改）：p.19 “innder product”（应为inner product）、p.19 “inner produce”（应为inner product）、p.22 “hyeprsurfaces”（应为hypersurfaces）、p.22 “the the ζ space”（多一个 “the”）；另 p.18 “Fayet-Illiopoulos”（应为Fayet–Iliopoulos，见 tr02a 的批注）。这些都不影响数学内容，故只在此说明。]
